\section{Introdução e Pivot Estratégico}

Esta investigação transcende a visão binária e estática da falência empresarial, adotando uma perspetiva dinâmica e recorrente. O fenómeno central é o \textit{Recurring Financial Distress}, analisado através do ciclo \textbf{saúde-doença-recaída}.

\subsection{Diferenciação Crítica: Estados Transientes vs. Absorventes}
A literatura tradicional colapsa frequentemente diferentes tipos de eventos numa única variável de "distress". Para atingir a excelência preditiva, distinguimos três estados fundamentais:
\begin{itemize}
    \item \textbf{Estado 0 (Saudável):} Operação normal sem sinais de alerta.
    \item \textbf{Estado 1 (Distress Transiente):} Verificação de critérios operacionais recuperáveis (ex: Capital Próprio negativo ou EBITDA insuficiente para cobrir juros).
    \item \textbf{Estado 2 (Falência/Terminal):} Verificação de critérios de dissolução legal ou liquidação (Estado Absorvente).
\end{itemize}

\section{Hipóteses de Investigação}

\begin{description}
    \item[H1 (Feature Value):] A inclusão de variáveis de \textit{Corporate Governance} e Macroeconomia reduz significativamente o erro de calibração (Brier Score) e melhora a discriminação (C-Index) em comparação com modelos puramente financeiros baseados em rácios do SABI.
    \item[H2 (Superioridade Arquitetural):] Modelos de Deep Learning Sequencial (\textbf{DynamicDeepHit}) superam modelos estáticos (Cox, RSF) ao capturarem a memória financeira das PME através de trajetórias longitudinais.
    \item[H3 (Dinâmica da Recorrência):] Os determinantes da recaída (transição $0 \to 1$ após um episódio prévio) são estruturalmente diferentes dos determinantes do primeiro episódio de stress, justificando o uso de modelos estratificados (\textbf{Cox PWP}).
\end{description}

\section{Inovações Técnicas Centrais}
\begin{enumerate}
    \item \textbf{DynamicDeepHit:} Arquitetura com RNN/LSTM para processar sequências temporais e riscos concorrentes.
    \item \textbf{Cox PWP (Prentice-Williams-Peterson):} Modelação da recorrência com estratificação por episódio.
    \item \textbf{Calibração de Precisão:} Foco no \textit{Integrated Brier Score} (IBS) para medir a qualidade das probabilidades estimadas.
\end{enumerate}
